\documentclass[12pt]{article}
\usepackage[utf8]{inputenc}
\usepackage[T1]{fontenc}
\usepackage[margin=1in]{geometry}
\usepackage{setspace}
\setlength{\parindent}{2em}
\setlength{\parskip}{1.25em}
\renewcommand{\baselinestretch}{1.0}
\makeatletter
\newcommand{\@subtitle}{}
\newcommand{\subtitle}[1]{\renewcommand{\@subtitle}{#1}}
\renewcommand{\@maketitle}{%
  \newpage
  \null
  \vskip 1em%
  \begin{center}%
    {\LARGE \@title \par}%
    \ifx\@subtitle\@empty\else
      \vskip 0.5em%
      {\large \@subtitle \par}%
    \fi
  \end{center}%
  \par
  \vskip 1em}
\makeatother
\usepackage{xcolor}
\usepackage[colorlinks=true,urlcolor=blue!55!black,linkcolor=black]{hyperref}

\title{\textit{Calder v. Bull}, 3 U.S. (3 Dall.) 386 (1798)}
\subtitle{Case Brief}
\author{}
\date{}

\begin{document}
\maketitle

\section{Memory Jogger}

In \href{https://www.courtlistener.com/opinion/84695/calder-v-bull/}{\textit{Calder v. Bull}}, the Supreme Court unanimously held that the Constitution's Ex Post Facto Clause prohibits only retroactive criminal laws---not retroactive civil legislation---while sparking the foundational American debate between Justice Chase's natural-law constitutionalism and Justice Iredell's textual positivism.

\section{Detailed Case Facts}

On August 21, 1779, Normand Morrison (the grandson) executed a will leaving property to Caleb Bull and his wife (\href{https://supreme.justia.com/cases/federal/us/3/386/}{3 U.S. at 386--387}). On March 21, 1793, the Court of Probate for Hartford County, Connecticut, issued a decree disapproving of Normand Morrison's will and refusing to record it (\href{https://supreme.justia.com/cases/federal/us/3/386/}{3 U.S. at 386}). Because the will was disapproved, the estate passed by intestacy to Calder and his wife, who claimed the property in right of Mrs. Calder as heiress of Normand Morrison, physician (\href{https://supreme.justia.com/cases/federal/us/3/386/}{3 U.S. at 387}).

Under an existing Connecticut statute, parties aggrieved by a probate decree had 18 months to appeal. More than 18 months elapsed after the March 1793 decree without an appeal by Caleb Bull and his wife, extinguishing their statutory right of appeal, and no general Connecticut statute authorized a new hearing before the probate court (\href{https://supreme.justia.com/cases/federal/us/3/386/}{3 U.S. at 386--387}).

In 1795, however, Connecticut did not yet have a modern written state constitution separating legislative and judicial powers; its government operated under its colonial charter, acts of assembly, and long-standing usages under which the Connecticut Legislature (the General Court) historically exercised superintending judicial authority, including granting new trials (\href{https://press-pubs.uchicago.edu/founders/documents/a1_10_1s10.html}{3 U.S. at 392--393, 395--396}). On the second Thursday of May 1795, the Connecticut Legislature passed a special resolution or law setting aside the Hartford Court of Probate's March 21, 1793 decree and granting a new hearing before the same probate court, with liberty to appeal within six months (\href{https://supreme.justia.com/cases/federal/us/3/386/}{3 U.S. at 386}).

At the new hearing held pursuant to the legislative resolution, the Hartford Court of Probate reversed its earlier ruling and, on July 27, 1795, approved Normand Morrison's will and ordered it recorded, thereby awarding the property to Bull and his wife (\href{https://supreme.justia.com/cases/federal/us/3/386/}{3 U.S. at 386--387}). Calder and his wife challenged the Connecticut Legislature's 1795 resolution as an unconstitutional \textit{ex post facto} law barred by Article I, \S 10, cl. 1 of the U.S. Constitution, which provides that ``No State shall ... pass any ... ex post facto Law.''

\section{Procedural History}

Following the Hartford Court of Probate's July 27, 1795 decree approving the will, Calder and his wife appealed in August 1795 to the Superior Court at Hartford, which affirmed the probate decree at its February 1796 term (\href{https://supreme.justia.com/cases/federal/us/3/386/}{3 U.S. at 386}). Calder and his wife then appealed to the Supreme Court of Errors of Connecticut, the state's highest court, which in June 1796 adjudged that there were no errors and upheld the legislative resolution under both Connecticut law and the federal Constitution (\href{https://supreme.justia.com/cases/federal/us/3/386/}{3 U.S. at 386, 393}).

Calder and his wife brought the case to the U.S. Supreme Court on writ of error from the State of Connecticut. The case was argued on February 8 and February 13, 1798, and decided on August 8, 1798 (\href{https://www.oxfordreference.com/display/10.1093/oi/authority.20110803095542467}{Oxford Reference}; \href{https://supreme.justia.com/cases/federal/us/3/386/}{Justia}). The Supreme Court unanimously affirmed the judgment of the Supreme Court of Errors of Connecticut (\href{https://supreme.justia.com/cases/federal/us/3/386/}{3 U.S. at 401}).

\section{Judicial Votes}

Following the pre-Marshall Court practice, the four participating Justices delivered individual \textit{seriatim} opinions rather than a single opinion of the Court (\href{https://www.oxfordreference.com/display/10.1093/oi/authority.20110803095542467}{Oxford Reference}).

\begin{itemize}
  \item \textbf{Unanimous Judgment Affirming (4--0, Seriatim Opinions):}
  \begin{itemize}
    \item Justice Samuel Chase (3 U.S. at 386--395)
    \item Justice William Paterson (3 U.S. at 395--397)
    \item Justice James Iredell (3 U.S. at 398--400)
    \item Justice William Cushing (3 U.S. at 400--401)
  \end{itemize}
  \item \textbf{Concurrences:} None separate from the four \textit{seriatim} opinions (all four Justices concurred in the judgment affirming Connecticut's judgment, though Justice Iredell expressly disagreed with Justice Chase's natural-law reasoning).
  \item \textbf{Dissents:} None.
  \item \textbf{Not Participating:} Chief Justice Oliver Ellsworth and Justice James Wilson did not participate in the decision (\href{https://www.oxfordreference.com/display/10.1093/oi/authority.20110803095542467}{Oxford Reference}).
\end{itemize}

\section{Holding}

\textbf{Judgment:} Affirmed. The Connecticut Legislature's May 1795 resolution setting aside a probate decree and granting a new hearing did not violate the federal Constitution (\href{https://supreme.justia.com/cases/federal/us/3/386/}{3 U.S. at 386, 401}).

\textbf{Rules Announced:}
\begin{enumerate}
  \item \textbf{Criminal-Only Scope of the Ex Post Facto Clause:} The constitutional prohibition on \textit{ex post facto} laws (Art. I, \S\S 9 and 10) applies solely to retroactive criminal (penal) laws, not to retroactive civil legislation affecting private property or judicial proceedings. Justice Chase's canonical four-category definition encompasses only laws that: (1) criminalize and punish an action that was innocent when done; (2) aggravate a crime or make it greater than when committed; (3) inflict a greater punishment than annexed to the crime when committed; or (4) alter the legal rules of evidence to permit conviction on less or different testimony than required at the time of the offense (\href{https://supreme.justia.com/cases/federal/us/3/386/}{3 U.S. at 390}). Laws that ``mollify the rigor of the criminal law'' are not prohibited (\href{https://supreme.justia.com/cases/federal/us/3/386/}{3 U.S. at 391}).
  \item \textbf{No Federal Jurisdiction Over Pure State Constitutional Claims:} The U.S. Supreme Court has no jurisdiction to declare a state statute void on the ground that it violates the state's own constitution (\href{https://supreme.justia.com/cases/federal/us/3/386/}{3 U.S. at 392}).
  \item \textbf{Deference to State Allocation of Judicial Power:} Where a state's own constitutional usage permits its legislature to act in a judicial capacity by awarding new trials, the federal Constitution does not bar that practice (\href{https://supreme.justia.com/cases/federal/us/3/386/}{3 U.S. at 395--396, 398, 400--401}).
\end{enumerate}

\textbf{Holding vs. Dicta:} The narrow holding is that Connecticut's 1795 resolution granting a new probate hearing was not a prohibited \textit{ex post facto} law under Article I, \S 10 because the Clause governs only criminal laws (and, for Paterson, Iredell, and Cushing, because the resolution could also be viewed as a customary judicial act). The famous philosophical clash between Justice Chase---who asserted that acts contrary to ``the great first principles of the social compact'' are void even without an explicit constitutional prohibition---and Justice Iredell---who insisted courts may invalidate laws only when they violate the written Constitution---was pure \textbf{dicta}, because both Justices voted to uphold Connecticut's resolution.

\section{Analysis of Opinions}

\subsection{Justice Chase (Seriatim)}

Justice Chase structured his opinion around two distinct themes: first, a bold natural-law theory of limited republican government, and second, a rigorous textual and historical narrowing of the Ex Post Facto Clause. Opening with social-compact theory, Chase declared that he could not ``subscribe to the omnipotence of a state legislature'' even where a state constitution lacked express restraints (\href{https://supreme.justia.com/cases/federal/us/3/386/}{3 U.S. at 387--388}). Certain ``vital principles in our free republican governments''---such as prohibiting a legislature from making a man a judge in his own cause, impairing lawful private contracts, or taking property from A and giving it to B---limit legislative authority because ``it cannot be presumed'' that the people ever entrusted a legislature with such powers (\href{https://supreme.justia.com/cases/federal/us/3/386/}{3 U.S. at 388--389}).

Turning to the federal Ex Post Facto Clause, however, Chase adopted a disciplined historical and structural construction. Drawing on English Parliamentary history (such as the attainders of the Earl of Strafford and Sir John Fenwick), Blackstone's \textit{Commentaries}, Wooddeson, \textit{The Federalist}, and the declarations of rights of Massachusetts, Maryland, North Carolina, and Delaware, Chase concluded that ``ex post facto laws'' was a technical legal phrase referring exclusively to retroactive penal statutes (\href{https://supreme.justia.com/cases/federal/us/3/386/}{3 U.S. at 389--392}). Structurally, if the Ex Post Facto Clause barred all retroactive civil laws, then Article I, \S 10's separate prohibitions on laws impairing the obligation of contracts and non-gold/silver legal tender---as well as the Fifth Amendment's Takings Clause---would have been ``unnecessary, and therefore improper'' (\href{https://supreme.justia.com/cases/federal/us/3/386/}{3 U.S. at 390, 393--394}). Chase also added that the Supreme Court had no jurisdiction to enforce state constitutions against state legislatures, and that the 1793 probate decree had not vested an exclusive property right in Calder anyway (\href{https://supreme.justia.com/cases/federal/us/3/386/}{3 U.S. at 392, 394}).

\subsection{Justice Paterson (Seriatim)}

Justice Paterson offered two independent grounds to affirm. First, because Connecticut's unwritten constitution rested on historical usage, the General Court acted in a ``double capacity'' as both legislature and court of judicature; awarding a new trial could therefore be respected as a customary judicial act rather than a legislative one (\href{https://supreme.justia.com/cases/federal/us/3/386/}{3 U.S. at 395--396}). Second, even treating the resolution as legislative, Paterson agreed with Chase that ``ex post facto'' had an established technical meaning limited to ``crimes, pains, and penalties,'' confirmed by Blackstone, early state constitutions, and the text of Article I, \S 10's neighboring Contracts Clause (\href{https://supreme.justia.com/cases/federal/us/3/386/}{3 U.S. at 396--397}). Remarkably, Paterson candidly admitted that he had felt ``an ardent desire to have extended the provision in the Constitution to retrospective laws in general''---which he personally disliked---but was compelled by the text and technical meaning to confine the Clause to criminal laws (\href{https://supreme.justia.com/cases/federal/us/3/386/}{3 U.S. at 397}).

\subsection{Justice Iredell (Seriatim)}

Justice Iredell concurred in the result---agreeing both that Connecticut's legislature exercised a traditional judicial power analogous to England's House of Lords and that the Ex Post Facto Clause applies only to criminal cases---but delivered a celebrated rebuttal to Justice Chase's natural-law dicta (\href{https://supreme.justia.com/cases/federal/us/3/386/}{3 U.S. at 398--400}). Iredell argued that American constitutions were written precisely to define legislative boundaries with clarity. If a legislature violates a written constitutional provision, courts must declare the act void (though, given the ``delicate and awful nature'' of that power, only in a ``clear and urgent case'') (\href{https://supreme.justia.com/cases/federal/us/3/386/}{3 U.S. at 399}; \href{https://www.heritage.org/courts/report/judicial-supremacy-and-our-two-constitutions-reflections-the-historical-record}{Heritage Foundation}). But if a legislature acts within its delegated constitutional powers, courts cannot strike the law down merely because judges believe it violates ``natural justice,'' for ``[t]he ideas of natural justice are regulated by no fixed standard; the ablest and the purest men have differed upon the subject'' (\href{https://supreme.justia.com/cases/federal/us/3/386/}{3 U.S. at 399}). Iredell further noted that civil retroactivity is sometimes essential to public governance, such as exercising eminent domain for highways and fortifications with reasonable compensation (\href{https://supreme.justia.com/cases/federal/us/3/386/}{3 U.S. at 400}).

\subsection{Justice Cushing (Seriatim)}

Justice Cushing wrote a two-sentence \textit{seriatim} opinion finding the case ``clear of all difficulty, taken either way'': if Connecticut's resolution was a judicial act, the federal Constitution did not touch it; if it was a legislative act, it was ``maintained and justified by the ancient and uniform practice of the State of Connecticut'' (\href{https://supreme.justia.com/cases/federal/us/3/386/}{3 U.S. at 400--401}).

\section{Examples: Future Applications}

\begin{enumerate}
  \item \textbf{Same-Side Application 1 (Retroactive Civil Tax or Liability Statute):} Congress enacts a retroactive amendment closing an estate-tax loophole for transactions completed during the preceding calendar year. Under \textit{Calder v. Bull}, the amendment does \textbf{not} violate the Ex Post Facto Clause because it is a civil revenue measure rather than a criminal punishment, even though it impairs settled financial expectations (though it remains subject to Fifth Amendment Due Process review).
  \item \textbf{Same-Side Application 2 (Retroactive Revocation of Civil Professional License):} A state passes a statute barring any person with a prior felony conviction from holding a medical license or waterfront union office, applying it to individuals convicted years earlier. Under \textit{Calder}, because the measure is a nonpunitive civil regulation protecting public safety rather than an added criminal penalty, it is \textbf{not} an unconstitutional \textit{ex post facto} law (\href{https://law.justia.com/constitution/us/article-1/70-ex-post-facto-laws.html}{U.S. Constitution Annotated}).
  \item \textbf{Opposite-Side Application 1 (Retroactive Reduction of Evidence Required to Convict):} A state amends its criminal code to eliminate the requirement that a complainant's testimony be corroborated in certain sexual-offense prosecutions, and applies the change to offenses committed before the amendment's effective date. Under \textit{Calder}'s fourth category---laws altering ``the legal rules of evidence'' to allow conviction on ``less or different testimony''---this law is an \textbf{unconstitutional} \textit{ex post facto} law (see \textit{Carmell v. Texas}, 529 U.S. 513 (2000)).
  \item \textbf{Opposite-Side Application 2 (Retroactive Mandatory Criminal Restitution or Sentence Enhancement):} Congress increases the mandatory prison term or extends the enforceable period of a punitive criminal restitution order for defendants whose offenses occurred before the new statute took effect. Under \textit{Calder}'s third category---laws inflicting a ``greater punishment than the law annexed to the crime, when committed''---retroactive application is \textbf{unconstitutional} (see \href{https://www.supremecourt.gov/opinions/25pdf/24-482_d1oe.pdf}{\textit{Ellingburg v. United States} (2026)}).
  \item \textbf{Fence-Sitter (Retroactive Sex-Offender Registration and Residency Restrictions):} A state enacts a retroactive sex-offender statute labeled ``civil and regulatory'' that mandates lifetime public internet registration, in-person reporting, and strict geographic residency bans for offenders convicted before the law existed. Under \textit{Calder}, if the scheme is genuinely civil and nonpunitive, the Ex Post Facto Clause does not apply (see \textit{Smith v. Doe}, 538 U.S. 84 (2003)); however, if the punitive effects (such as banishment-like residency bars or coercive penalties for a public wrong) override the legislature's civil label, it crosses the line into prohibited retroactive punishment (\href{https://www.supremecourt.gov/opinions/25pdf/24-482_d1oe.pdf}{\textit{Ellingburg v. United States} (2026)}).
\end{enumerate}

\section{Critique}

\textit{Calder v. Bull} occupies an enduring and controversial place in American constitutional law on two fronts: its narrowing of the Ex Post Facto Clause to criminal laws and its iconic debate over unenumerated natural-law limits on legislation. Logically, commentators have noted an internal tension in Justice Chase's opinion: after passionately declaring that a law taking property from A and giving it to B violates ``the great first principles of the social compact,'' Chase declined to invalidate a Connecticut resolution that functionally shifted an inheritance from Calder to Bull, resting on the technicality that the 1793 probate decree had not ``vested'' a perfect property right in Calder.

\textbf{Progressive Critique:} Progressive and modern living-constitutionalist scholars view the Chase--Iredell debate with mixed admiration and caution. On one hand, many progressives embrace the spirit of Justice Chase's refusal to accept legislative omnipotence as an early precursor to substantive due process and the protection of fundamental unenumerated rights (such as privacy and bodily autonomy). On the other hand, progressives warn that Chase's specific examples---protecting vested property and private contracts against legislative redistribution---foreshadowed the \textit{Lochner} era's judicial hostility to economic and social-welfare regulation, and they generally support \textit{Calder}'s actual holding leaving legislatures broad leeway to enact retroactive civil, tax, and remedial economic legislation (see \href{https://kb.osu.edu/bitstreams/177d2af3-9d00-5fe4-af96-0073e9b3e460/download}{Ohio State Law Journal}).

\textbf{Originalist and Textualist Critique:} Originalists and textualists are divided on \textit{Calder}'s two legacies. Nearly all modern originalists celebrate Justice Iredell's positivist rebuttal to Chase as the canonical statement of written constitutionalism: judges derive their authority to strike down laws solely from the ratified text of the Constitution, never from abstract, judge-made ``natural justice'' (\href{https://www.heritage.org/courts/report/judicial-supremacy-and-our-two-constitutions-reflections-the-historical-record}{Heritage Foundation}). Yet prominent originalists have sharply questioned \textit{Calder}'s substantive Ex Post Facto holding. In \href{https://www.law.cornell.edu/supct/html/97-42.ZC.html}{\textit{Eastern Enterprises v. Apfel}, 524 U.S. 498, 538--539 (1998) (Thomas, J., concurring)}, Justice Clarence Thomas expressed skepticism of \textit{Calder}'s criminal-only limitation---noting that Chase's redundancy argument regarding the Takings Clause was unpersuasive---and signaled willingness to reconsider whether the Ex Post Facto Clauses originally reached retroactive civil laws. More recently, in \href{https://www.supremecourt.gov/opinions/25pdf/24-482_d1oe.pdf}{\textit{Ellingburg v. United States} (2026)}, originalist analysis has similarly argued that in 1798 ``punishment'' for a ``crime'' encompassed coercive penalties for public wrongs that modern courts too easily exempt by deferring to a legislature's ``civil'' label.

\section{Key Quotations}

\begin{quote}
An act of the legislature (for I cannot call it a law) contrary to the great first principles of the social compact cannot be considered a rightful exercise of legislative authority.
\end{quote}
\noindent --- Justice Chase, \href{https://supreme.justia.com/cases/federal/us/3/386/}{3 U.S. at 388}.

\begin{quote}
1st. Every law that makes an action done before the passing of the law and which was innocent when done, criminal and punishes such action. 2d. Every law that aggravates a crime, or makes it greater than it was when committed. 3rd. Every law that changes the punishment, and inflicts a greater punishment than the law annexed to the crime, when committed. 4th. Every law that alters the legal rules of evidence and receives less or different testimony than the law required at the time of the commission of the offense in order to convict the offender.
\end{quote}
\noindent --- Justice Chase, \href{https://supreme.justia.com/cases/federal/us/3/386/}{3 U.S. at 390}.

\begin{quote}
Every ex post facto law must necessarily be retrospective, but every retrospective law is not an ex post facto law. The former only are prohibited.
\end{quote}
\noindent --- Justice Chase, \href{https://supreme.justia.com/cases/federal/us/3/386/}{3 U.S. at 391}.

\begin{quote}
The words ``ex post facto,'' when applied to a law, have a technical meaning, and in legal phraseology refer to crimes, pains, and penalties.
\end{quote}
\noindent --- Justice Paterson, \href{https://supreme.justia.com/cases/federal/us/3/386/}{3 U.S. at 396}.

\begin{quote}
If, on the other hand, the legislature of the Union or the legislature of any member of the Union shall pass a law within the general scope of its constitutional power, the court cannot pronounce it to be void merely because it is in its judgment contrary to the principles of natural justice. The ideas of natural justice are regulated by no fixed standard; the ablest and the purest men have differed upon the subject, and all that the court could properly say in such an event would be that the legislature (possessed of an equal right of opinion) had passed an act which, in the opinion of the judges, was inconsistent with the abstract principles of natural justice.
\end{quote}
\noindent --- Justice Iredell, \href{https://supreme.justia.com/cases/federal/us/3/386/}{3 U.S. at 399}.

\end{document}
